% TOP_PROBLEM: {"id": 30006722, "title": "Grothendieck–Teichmüller Conjecture", "category_id": 1, "category": {"id": 1, "name": "number_theory", "display_name": "Number Theory", "description": "Properties of integers, prime numbers, Diophantine equations.", "slug": "number-theory", "order_index": 1, "created_at": "2026-07-31T15:26:25.670Z"}, "status": "open"}
% FIRST_POSTED: 2026-09-25
% ENTRY_KIND: statement_only
% !TeX program = lualatex
% Definition and sources only; this file is not an LLM solution attempt.
\documentclass[11pt]{article}
\usepackage[margin=25mm]{geometry}
\usepackage{fontspec}
\setmainfont{Latin Modern Roman}
\usepackage{amsmath,amssymb,mathtools}
\usepackage{unicode-math}
\setmathfont{Latin Modern Math}
\usepackage{xurl,hyperref}
\hypersetup{colorlinks=true,urlcolor=blue}
\setcounter{secnumdepth}{0}
\setlength{\parindent}{0pt}
\setlength{\parskip}{5pt}
\setlength{\emergencystretch}{3em}
\begin{document}
\section*{\texorpdfstring{Grothendieck–Teichmüller Conjecture}{Grothendieck–Teichmüller Conjecture}}\label{30006722}
\subsection{Definitions and mathematical statement}
Let $\widehat F_2$ be the free profinite group on $x,y$, and let $\chi:G_{{\mathbb{Q}}}\to\widehat{{\mathbb{Z}}}^\times$ be the cyclotomic character. Drinfeld's original $\widehat{\mathrm{GT}}$ consists of the compatible invertible profinite pairs satisfying the commutator, symmetry, hexagon, and pentagon conditions; all these conditions, not only the pair equations, are incorporated from the source. Under the fixed Ihara splitting,
\[
 \mathrm{Ih}(g)=\bigl((\chi(g)-1)/2,f_g\bigr),\qquad
 x\mapsto x^{\chi(g)},\quad y\mapsto f_g^{-1}y^{\chi(g)}f_g.
\]
The question is surjectivity of this actual embedding $G_{{\mathbb{Q}}}\hookrightarrow\widehat{\mathrm{GT}}$. Gentle, pro-$\ell$, and prounipotent groups are not substituted. Composition and coordinate conventions must agree with Section 1.3 of the cited paper.
\par\smallskip\textit{Formulation and concept references:} \href{https://arxiv.org/pdf/2405.11725v3}{[S1]}.

\subsection{Short English statement}
Is every symmetry in the original profinite Grothendieck--Teichmüller group induced by a unique automorphism of the algebraic closure of the rationals?
\subsection{Sources}
\begin{itemize}
\item[S1]Ivan Bortnovskyi, Vasily A. Dolgushev, Borys Holikov, Vadym Pashkovskyi, Accessing non-abelian quotients of the Grothendieck-Teichmueller group via elementary tools, arXiv2405.11725v3,6August2026,35pages.
\url{https://arxiv.org/pdf/2405.11725v3}.
\textit{Formulation source. Consulted 24 September 2026: Section 1.3, equations (1.5)--(1.6).}
\item[S2]Accessing non-abelian quotients of the Grothendieck-Teichmueller group via elementary tools.
\url{https://arxiv.org/html/2405.11725}.
\textit{Further reference; not a proof endorsement.}
\item[S3]Galois actions on surfaces and a higher genus Grothendieck-Teichmüller group.
\url{https://arxiv.org/html/2606.01466v1}.
\textit{Further reference; not a proof endorsement.}
\item[S4]Proving the Grothendieck–Teichmüller Conjecture for Profinite Spaces \& The Galois Grothendieck Path Integral.
\url{https://arxiv.org/html/2503.13006}.
\textit{Further reference; not a proof endorsement.}
\end{itemize}
\end{document}
